% GENERATED FILE. Edit canonical lessons, then run npm run generate:books.
% canonical-source-hash: fnv1a64:2c07432c45bce8c6
% canonical-lessons: ES-C713-diversion, ES-C713-divertido, ES-C713-gracioso, ES-C713-aburrirse, ES-C713-aburrimiento

\chapter{Pasarlo bien}
\label{ch:es-pasarlo-bien}
\hlchaptermodality{713}

\begin{chapteropening}
\textbf{By the end of this chapter:} \emph{I can say whether something was fun, funny or boring.}

Complete the last lesson of chapter 713: I can say whether something was fun, funny or boring.
\end{chapteropening}

\section[la diversión]{la diversión — fun, entertainment}

\label{lesson:ES-C713-diversion}



\begin{quote}
Before the new one: say the Spanish for to cheer up, then the Spanish for relief.
\end{quote}

\subsection*{You'll want to know: la diversión}
\textbf{la diversión} — \textquotedblleft{}fun, entertainment\textquotedblright{}.

\begin{quote}
Para los niños, la playa es pura diversión.
\end{quote}

\begin{quote}
\emph{For children, the beach is pure fun.}
\end{quote}

\emph{¡Qué diversión!} — What fun!

\begin{cousinweb}
\emph{diversión} is built on \emph{divertir}, from Latin \emph{dīvertere} \textquotedblleft{}to turn aside\textquotedblright{} (EtymDB), from \emph{vertere} \textquotedblleft{}to turn\textquotedblright{}: fun is what turns you away from your cares.
\end{cousinweb}

\begin{grammarlens}[title={What you've built}]
One more word for telling, explaining and describing at B1.
\end{grammarlens}

\subsection*{Your turn}
\begin{itemize}
\raggedright
  \item \emph{Say it:} \emph{la diversión}
  \item \emph{Say it:} \emph{la diversión}, once more
  \item \emph{Say it:} a sentence of your own with \emph{la diversión}
  \item \emph{Recall:} say the Spanish for to cheer up, then the Spanish for relief, then say \emph{la diversión} again
\end{itemize}

\begin{tcolorbox}[breakable,colback=teal!4,colframe=teal!35!black,title={Before you move on}]
What does la diversión mean? (\textquotedblleft{}Fun, entertainment\textquotedblright{}.) Say it once more.
\end{tcolorbox}
\section[divertido]{divertido — entertaining, fun}

\label{lesson:ES-C713-divertido}



\begin{quote}
Before the new one: say the Spanish for relief, then the Spanish for fun.
\end{quote}

\subsection*{You'll want to know: divertido}
\textbf{divertido} — \textquotedblleft{}entertaining, fun\textquotedblright{}.

\begin{quote}
La fiesta fue muy divertida.
\end{quote}

\begin{quote}
\emph{The party was a lot of fun.}
\end{quote}

With \emph{ser} it means fun; \emph{divertirse} is to have a good time.

\begin{cousinweb}
\emph{divertido} is the participle of \emph{divertir}, which comes from Latin \emph{dīvertere} \textquotedblleft{}to turn aside\textquotedblright{} (EtymDB), built on \emph{vertere} \textquotedblleft{}to turn\textquotedblright{}.
\end{cousinweb}

\begin{grammarlens}[title={What you've built}]
One more word for telling, explaining and describing at B1.
\end{grammarlens}

\subsection*{Your turn}
\begin{itemize}
\raggedright
  \item \emph{Say it:} \emph{divertido}
  \item \emph{Say it:} \emph{divertido}, once more
  \item \emph{Say it:} a sentence of your own with \emph{divertido}
  \item \emph{Recall:} say the Spanish for relief, then the Spanish for fun, then say \emph{divertido} again
\end{itemize}

\begin{tcolorbox}[breakable,colback=teal!4,colframe=teal!35!black,title={Before you move on}]
What does divertido mean? (\textquotedblleft{}Entertaining, fun\textquotedblright{}.) Say it once more.
\end{tcolorbox}
\section[gracioso]{gracioso — funny, amusing}

\label{lesson:ES-C713-gracioso}



\begin{quote}
Before the new one: say the Spanish for fun, then the Spanish for entertaining.
\end{quote}

\subsection*{You'll want to know: gracioso}
\textbf{gracioso} — \textquotedblleft{}funny, amusing\textquotedblright{}.

\begin{quote}
Tu hermano es muy gracioso.
\end{quote}

\begin{quote}
\emph{Your brother is very funny.}
\end{quote}

\emph{Gracioso} is funny in a pleasant way; \emph{chistoso} is close.

\begin{cousinweb}
\emph{gracioso} was borrowed from Latin \emph{grātiōsus} \textquotedblleft{}agreeable, popular\textquotedblright{} (EtymDB), built on \emph{grātia} \textquotedblleft{}favour, charm\textquotedblright{}, the source of \emph{gracia}.
\end{cousinweb}

\begin{grammarlens}[title={What you've built}]
One more word for telling, explaining and describing at B1.
\end{grammarlens}

\subsection*{Your turn}
\begin{itemize}
\raggedright
  \item \emph{Say it:} \emph{gracioso}
  \item \emph{Say it:} \emph{gracioso}, once more
  \item \emph{Say it:} a sentence of your own with \emph{gracioso}
  \item \emph{Recall:} say the Spanish for fun, then the Spanish for entertaining, then say \emph{gracioso} again
\end{itemize}

\begin{tcolorbox}[breakable,colback=teal!4,colframe=teal!35!black,title={Before you move on}]
What does gracioso mean? (\textquotedblleft{}Funny, amusing\textquotedblright{}.) Say it once more.
\end{tcolorbox}
\section[aburrirse]{aburrirse — to get bored}

\label{lesson:ES-C713-aburrirse}



\begin{quote}
Before the new one: say the Spanish for entertaining, then the Spanish for funny.
\end{quote}

\subsection*{You'll want to know: aburrirse}
\textbf{aburrirse} — \textquotedblleft{}to get bored\textquotedblright{}.

\begin{quote}
Me aburro en las reuniones largas.
\end{quote}

\begin{quote}
\emph{I get bored in long meetings.}
\end{quote}

Reflexive: \emph{me aburro}. \emph{Aburrir} without \emph{se} is to bore someone.

\begin{cousinweb}
\emph{aburrirse} is inherited from Latin \emph{abhorrēre} \textquotedblleft{}to shrink from in horror\textquotedblright{}, through a spoken form \emph{abhorrīre} (EtymDB); English \emph{abhor} is the same word.
\end{cousinweb}

\begin{grammarlens}[title={What you've built}]
One more word for telling, explaining and describing at B1.
\end{grammarlens}

\subsection*{Your turn}
\begin{itemize}
\raggedright
  \item \emph{Say it:} \emph{aburrirse}
  \item \emph{Say it:} \emph{aburrirse}, once more
  \item \emph{Say it:} a sentence of your own with \emph{aburrirse}
  \item \emph{Recall:} say the Spanish for entertaining, then the Spanish for funny, then say \emph{aburrirse} again
\end{itemize}

\begin{tcolorbox}[breakable,colback=teal!4,colframe=teal!35!black,title={Before you move on}]
What does aburrirse mean? (\textquotedblleft{}To get bored\textquotedblright{}.) Say it once more.
\end{tcolorbox}
\section[el aburrimiento]{el aburrimiento — boredom}

\label{lesson:ES-C713-aburrimiento}



\begin{quote}
Before the new one: say the Spanish for funny, then the Spanish for to get bored.
\end{quote}

\subsection*{You'll want to know: el aburrimiento}
\textbf{el aburrimiento} — \textquotedblleft{}boredom\textquotedblright{}.

\begin{quote}
¡Qué aburrimiento de película!
\end{quote}

\begin{quote}
\emph{What a boring film!}
\end{quote}

\emph{¡Qué aburrimiento!} — How boring!

\begin{cousinweb}
\emph{aburrimiento} is built on \emph{aburrir} with the ending \emph{-miento} (EtymDB); the verb comes from Latin \emph{abhorrēre} \textquotedblleft{}to shrink from\textquotedblright{}.
\end{cousinweb}

\begin{grammarlens}[title={What you've built}]
One more word for telling, explaining and describing at B1.
\end{grammarlens}

\subsection*{Your turn}
\begin{itemize}
\raggedright
  \item \emph{Say it:} \emph{el aburrimiento}
  \item \emph{Say it:} \emph{el aburrimiento}, once more
  \item \emph{Say it:} a sentence of your own with \emph{el aburrimiento}
  \item \emph{Recall:} say the Spanish for funny, then the Spanish for to get bored, then say \emph{el aburrimiento} again
\end{itemize}

\begin{tcolorbox}[breakable,colback=teal!4,colframe=teal!35!black,title={Before you move on}]
What does el aburrimiento mean? (\textquotedblleft{}Boredom\textquotedblright{}.) Say it once more.
\end{tcolorbox}
